% !PRINTED 60
\weberpage{60}

\begin{center}
\S.\,16.

\gesperrt{The Groups of Number Classes with Respect to a Composite Modulus}.
\end{center}

From the combination of the two theorems 2. and 3. of the two preceding sections there now results the following result, by which the problem that we posed at the beginning of \S.\,14 is completely solved:

\noindent\textbf{4.} \gesperrt{If the modulus}

\[
m = 2^\lambda q_1^{\varkappa_1} q_2^{\varkappa_2} \ldots
\]

\gesperrt{is, and} $g_1, g_2 \ldots$ \gesperrt{are primitive roots of the squares of the primes} $q_1, q_2 \ldots$\gesperrt{, and if further}

\[
a = 2, \quad b = \tfrac12 \varphi(2^\lambda), \quad c_1 = \varphi(q_1^{\varkappa_1}), \quad c_2 = \varphi(q_2^{\varkappa_2}) \ldots
\]

\gesperrt{is set, then one can, for every number} $n$ \gesperrt{that is relatively prime to} $m$\gesperrt{, determine uniquely a system of numbers} $\alpha, \beta, \gamma_1, \gamma_2 \ldots$ \gesperrt{with respect to the moduli} $a, b, c_1, c_2 \ldots$\gesperrt{, such that the equations}

\weq{1}{
\begin{aligned}
n &\equiv (-1)^\alpha 5^\beta \pmod{2^\lambda} \\
&\equiv g_1^{\gamma_1} \pmod{q_1^{\varkappa_1}} \\
&\equiv g_2^{\gamma_2} \pmod{q_2^{\varkappa_2}} \\
&\; \cdots\cdots\cdots\cdots\cdots\cdots
\end{aligned}
}

\gesperrt{are satisfied.}

\gesperrt{The numbers $\alpha, \beta, \gamma_1, \gamma_2 \ldots$ are called the system of indices of $n$ for the modulus $m$.}

Here it is at first assumed that $\lambda \geqq 3$. But the theorem holds also for the remaining values of $\lambda$, if

\[
\text{for } \lambda = 0, 1, \quad a = 1, \quad b = 1
\]
\[
\text{\textquotedbl{}} \quad \lambda = 2, \qquad\quad a = 2, \quad b = 1
\]

is set. For $\lambda = 0, 1$ one is to set the indices $\alpha$ and $\beta = 0$, or else omit them entirely; for $\lambda = 2$, $\beta$ drops out, while $\alpha$ can be equal to 0 or equal to 1.

Now, in order to represent the group $\mathfrak{N}$ of the number classes $N$ with respect to the modulus $m$ by a basis, one determines the number classes

\weq{2}{A, B, C_1, C_2 \ldots,}

or at least representatives of these classes, which is always possible, from the congruences
%CONT>>
